\documentclass{article}
\usepackage[utf8]{inputenc}
\usepackage[english]{babel}
\usepackage{caption}
\usepackage{float}
\usepackage{subcaption}
\usepackage[table,xcdraw]{xcolor}

\usepackage[lmargin = {2.7cm},
rmargin = {2.7cm},
tmargin = {2.2cm},
bmargin = {2.2cm}
]{geometry}

\usepackage{comment}
\usepackage[hidelinks]{hyperref}
\usepackage[capitalise]{cleveref}

% TexStudio build config line: txs:///bibtex | txs:///biber | txs:///compile | txs:///view
\usepackage[
backend=biber,
style=numeric,
sorting=nty
]{biblatex}
\addbibresource{paper.bib}

\usepackage{graphicx}
\graphicspath{ {./} }

\title{Maschinenlesbare Ontologien und Vokabulare aus dem Bereich der industriellen Produktion}
\author{\textbf{Konrad Abicht} \\ k.abicht@gmail.com}
\date{TODO Datum}

\begin{document}

\maketitle

\paragraph{{\large Zusammenfassung}}

\paragraph{Einleitung:} Aktuell gibt es keine Liste von maschinenlesbaren, öffentlich-zugänglichen Ontologien und Vokabularen in RDF mit Bezug zum Themengebiet der industriellen Produktion. Ähnliche Arbeiten haben leidiglich eine kleine Menge an Ontologien zu dem Thema untersucht.

\paragraph{Methodik:} Es wurde im Zeitraum vom 01.11.2023 - 26.03.2024 eine Internet- und Literaturrecherche durchgeführt und eine CSV-Datei mit allen thematisch passenden bzw. referenzierten Ontologien und Vokabularen erstellt. Zu jeder Ontologie wurden eine Reihe von Metadaten erfasst (z.B. Titel, Kurzbeschreibung, Projektseite, Version, letzte Änderung und Lizenz).

\paragraph{Resultate:} Als Ergebnis wurden \textbf{343} Ontologien bzw. Vokabulare gefunden, wobei \textbf{189} einen thematischen Bezug zur industriellen Produktion haben. Nur 75 Ontologien besitzen vollständige Metadaten, bei dem Rest konnte mindestens ein Metadatum nicht ermittelt werden. Es wurden alle Ontologien ignoriert, zu denen es zwar wissenschaftliche Publikationen gibt, jedoch keine öffentlich einsehbaren RDF-Daten. Weitere Auswertungen sind in dieser Publikation zu finden. Die Rohdaten stehen unter einer freien Lizenz (CC-BY 4.0) zur freien Verwendung bereit.

\paragraph{Diskussion:} Die Untersuchung wurde lediglich über einen Zeitraum von 5 Monaten durchgeführt. Weiterhin wurde der thematische Fokus industrielle Produktion im weitesten Sinne betrachtet, wodurch möglicherweise auch thematisch nicht direkt passende Ontologien einbezogen wurden.

\paragraph{Fazit:} Diese Arbeit zeigt [...]

\newpage

\section{Einleitung}

Es wurde eine Internet- und Literaturrecherche zur Ermittlung aller öffentlich zugänglichen\footnote{RDF-Dokument ist über eine URL abrufbar.} und RDF vorliegenden Ontologien (und Vokabulare) aus dem Themengebiet der industriellen Produktion durchgeführt.
Das Ziel war die Erstellung einer Übersicht, bestehend aus maschinenlesbaren Ontologien, die eine automatisierte Verarbeitung des modellierten Wissens erlauben (z.B. OWL Reasoning oder Inferenz).

Die Arbeit ist wie folgt strukturiert: [...] Kapitel 2 enthält ... in Kapitel 3 folgt ... usw.

\subsection{Öffentlich verfügbare Forschungsdaten}

Die erstellten Forschungsdaten stehen in einem Github-Repository für die Öffentlichkeit zur Verfügung:

\textbf{TODO:} Link einfügen

\textbf{TODO:} Die Inhalte unterliegen dabei den Bedingungen der Creative Commons BY 4.0.

\subsubsection{Graph Visualisierung}

Skript: generate-yuml-diagram-config

- Vereinheitlichung der Prefixe bei den RDF-Daten (z.B. dc vs dcterms vs dcelements oder prov vs. prov\_namespace)

\section{Fachlicher Hintergrund}

In diesem Kapitel werden die notwendige Vorkenntnise für diese Arbeit kurz zusammengefasst.

\subsection{Semantic Web und formale Sprachen}

\textbf{TODO}
RDF/OWL

\subsection{Ontologien und Kontrollierte Vokabulare}

Im Rahmen dieser Arbeit lag der Fokus auf maschinenlesbaren RDF Ontologien und Vokabularen.
Eine Ontologie wurde aufgenommen, wenn sie die folgenden Kriterien erfüllt:

\begin{enumerate}
    \item Die Ontologie besitzt eine URL zu einer Textdatei, welche in einer formalen Sprache (RDF/OWL) vorliegt.
    \item Die Ontologie stellt ein Vokabular zur Beschreibung des Fachgebietes bereit.
    \item Der Ontologie liegt eine logische Theorie (z.B. in Form von Axiomen, Regeln, Hierarchien) über das Fachgebiet zugrunde, die auf das Vokabular zurückgreift.
\end{enumerate}

Diese Kriterien basieren größtenteils auf den Ausführungen von Fabian Neuhaus aus \cite{neuhaus2018ontology}, mit ein paar Ergänzungen.
Im Folgenden wird der Begriff Ontologie auch für Vokabulare genutzt, außer es wird explizit darauf hingewiesen.
Der Grund für diese Festlegung sind die Rechercheergebnisse, die zeigen, dass manche Autoren häufig ihre eigenen Arbeiten entweder als Ontologie oder Vokabular, oder beides, bezeichneten.
Zum Beispiel bezeichnet Martin Hepp GoodRelations als standardisiertes Vokabular für Produkte, Preise und Unternehmensdaten, nutzt jedoch Ontologie als Quasi-Synonym, siehe das folgende Zitat: "\textit{GoodRelations is a standardized vocabulary (also known as "schema", "data dictionary", or "ontology") for product, price, store, and company data that can [...]}"\footnote{Quelle: \url{https://www.heppnetz.de/ontologies/goodrelations/v1.html}}.

\subsubsection{Ontologiearten}

\begin{figure}[H]
    \includegraphics[scale=0.4]{ontology-types}
    \centering
    \caption{Ontology Types}
\end{figure}

Die Untersuchung ergab, dass die Autoren ihre Ontologien überwiegend nach der etablierten Vierteilung einordnen.
An der Spitze der Hierarchie steht die \textbf{Top-Level Ontologie} (Synonyme: Upper Ontology, Foundational Ontology).
Sie beinhaltet fachgebiets-unabhängige Inhalte (Begriffe, Relationen, Regeln, Axiome etc.) zur Beschreibung eines Ausschnitts der Realität.
Zu den bekanntesten gehören Suggested Upper Merged Ontology (SUMO), Descriptive Ontology for Linguistic and Cognitive Engineering (DOLCE) und Basic Formal Ontology (BFO).
Top-Level-Ontologien bringen häufig eine eigene Theorie zur Beschreibung der Realität mit.
Jede abgeleitete Ontologie übernimmt somit auch die zugehörige Theorie, Implikationen etc.

Es folgt die Ebene der \textbf{Mid-Level Ontologie}, welche die Inhalte der Top-Level Ontologie ergänzt bzw. verändert. Sie besitzen einen geringen thematischen Bezug und dient häufig als inhaltliche Brücke zwischen einer Top-Level Ontologie und ihr untergeordneten Kern- bzw. Task-Ontologien.

Danach folgt die Ebene der \textbf{Kern Ontologie} (Synonyme: Core Ontology, Domain Ontology).
Ihr Inhalt bezieht sich stärker auf ein Fachgebiet oder sie ergänzt die Inhalte anderer Ontologien.

Unter den Kern-Ontologien folgt die \textbf{Task-Ontologie} (Synonyme: Application Ontology).
Sie nutzt in der Regel die Inhalte von Ontologien höherer Ebenen und stellt eigene Inhalte für einen konkreten Anwendungsfall bereit.

\subsubsection{Ontology Design Pattern (ODP)}

Ontology Design Pattern (ODP) sind kleine unabhängige Ontologien, die einen sehr begrenzten Themenbereich modellieren und dabei einen hohen Grad an Wiederverwendbarkeit und Kombinierbarkeit mit anderen ODPs und Modellierungsansätzen besitzen \cite{gangemi2009ontology}.
Es besteht keine Beschränkung für deren Aufbau, Struktur und Ausrichtung.

\subsection{Industry Foundation Classes (IFC)}

Die Industry Foundation Classes (IFC) sind ein offener internationaler Standard zur digitalen Beschreibung von Gebäudemodellen im Building Information Model Format (kurz BIM)\footnote{\url
https://technical.buildingsmart.org/standards/ifc/}.
Es erlaubt den Datenasustausch zwischen Softwareanwendungen verschiedener Teilnehmer der Bau- und Facility-Management-Branche zur gemeinsamen Verwendung.
Der Standard enthält Definitionen für Daten, die für Gebäude während ihres Lebenszyklus erforderlich sind.
ICF wird an dieser Stelle kurz vorgestellt, weil es Ontologien mit IFC Bezug in den Rechercheergebnissen gibt.

\subsection{Lizenzinformationen}

Zu jeder Ontologie wurde die zugehörige Lizenz recherchiert.
Sie bestimmt die rechtlichen Rahmenbedingungen unter denen eine Ontologie benutzt werden darf.
Die verwendeten Lizenzen reichen von Public Domain bis hin zu restriktiven individuellen Lizenzregelungen.
Im Folgenden werden alle Lizenzen aufgelistet, die während der Recherche gefunden wurden:

\begin{enumerate}
    \item Apache License 2.0\footnote{\url{https://www.apache.org/licenses/LICENSE-2.0.html}}
    \item BSD License (2- und 3-Clause)\footnote{\url{https://opensource.org/license/bsd-2-clause}, \url{https://opensource.org/license/bsd-3-clause}}
    \item Creative Commons Licenses (verschiedene Arten)\footnote{\url{https://creativecommons.org/licenses/?lang=en}}
    \item GNU Public License (verschiedene Arten)\footnote{\url{https://www.gnu.org/licenses/licenses.en.html}}
    \item MIT License\footnote{\url{https://opensource.org/license/mit}}
    \item OGC Document License Agreement\footnote{\url{https://www.ogc.org/about-ogc/policies/document-license-agreement/}}
    \item Public Domain Dedication License (PDDL)\footnote{\url{https://opendatacommons.org/licenses/pddl/}}
    \item W3C Document License (2023 version)\footnote{\url{https://www.w3.org/copyright/document-license-2023/}}
\end{enumerate}


\section{Verwandte Arbeiten (Related Work)}

- Welche Untersuchungen wurden bereits dahingehend unternommen?

- Wozu gab es noch keine oder nur unzureichende Untersuchungen?

\subsection{Ontologies for Industry 4.0}
%
% Ontologies for Industry 4.0
% https://www.cambridge.org/core/services/aop-cambridge-core/content/view/BF86BB5310356D642C82470D67974804/S0269888919000109a.pdf/div-class-title-ontologies-for-industry-4-0-div.pdf

Zu Beginn der Publikation "Ontologies for Industry 4.0" \cite{kumar2019ontologies} (2019) geben die Autoren zuerst eine kurze Einordnung des Themengebietes Industrie 4.0, Factory 4.0 und Smart Manufacturing.
Diese wird ergänzt durch eine historische Einordnung.
Danach folgt eine Zusammenfassung relevanter Herausforderungen (z.B. Mensch-Maschine-Kommunikation oder Datenanalyse) und die zur Bewältigung genutzten Technologien (z.B. Internet of Things, 5G, Cloud).
Im zweiten Teil der Arbeit wird der ontologische Rahmen vorgestellt und die Autoren listen eine Reihe von Anwendungsbereichen für Ontologien auf, jedoch ohne diese näher einzuordnen.
Immerhin wird am Ende des Abschnitts ein kurzer Vergleich zwischen zwei verschiedenen Strukturierungsansätzen für Ontologien andiskutiert.
Im Kapitel 2.2 stellen die Autoren einige Resultate ihrer Standardisierungsvorhaben zu Industrieontologien vor und beschreiben den aktuellen Stand.
Das Hauptaugenmerk lag dabei zuerst auf der Kommunikation zwischen Robotern, menschlichen Bedienern, Endkunden und diversen Zulieferern.
Laut Aussage der Autoren sollten diese Themen zuerst geklärt werden, um daraus resultierende Probleme zu vermeiden.
Ihre Arbeit basiert auf dem IEEE-Beitrag namens "IEEE 1872-2015 Standard Ontologies for Robotics and Automation"\footnote{1872-2015 - IEEE Standard Ontologies for Robotics and Automation: \url{https://ieeexplore.ieee.org/document/7084073}}.
Diese Ontologien bieten eine ontologische Basis für weiterführende Konzepte aus dem Bereich Industrie 4.0.
Danach folgt eine kurze Vorstellung der folgenden Ontologien:
\begin{enumerate}
    \item \textbf{CORA:} Core Ontology for Robotics and Automation
    \item \textbf{ROA:} The Ontology for Autonomous Robotics\footnote{In der zugehörigen Publikation\cite{olszewska2017ontology} kürzen die Autoren die Ontologie jedoch mit ORA ab.}
    \item \textbf{ORArch}: Ontology for Robotic Architecture
    \item \textbf{O4I4}: Ontology for Industry 4.0
\end{enumerate}

Im dritten Teil der Arbeit gehen die Autoren die folgenden Szenarien im Bereich Industrie 4.0 ein: Smart-rapid prototyping scenario und UAV’s good delivery scenario.
Jedes Szenario wird kurz eingeführt und danach der Einsatz der Ontologien beschrieben.
Es ist erwähnenswert, dass die Publikation an mehreren Stellen eine Auflistung von relevanten Publikationen enthält, die jedoch nur benannt aber nicht inhaltlich eingeordnet werden. Damit bieten die Autoren einen guten Ausblick, jedoch fehlt dem Leser die Einordnung in das Papier.


\subsection{Where to Publish and Find Ontologies? A Survey of Ontology Libraries}

Die Autoren des Papiers "Where to Publish and Find Ontologies? A Survey of Ontology Libraries" \cite{d2012publish} (Jahr 2012) geben einen Überblick über Ontologie-Bibliotheken.
Es wurden verschiedene Bezeichnungen\footnote{Ontology Directory, Ontology Repository, Ontology Archive} für solche Systeme gefunden, weshalb man sich zum besseren Verständnis auf Ontologie-Bibliothek festgelegt hat.
Diese Publikation ist für diese Arbeit relevant, weil sie einen Überblick vorhandener Ontologie-Bibliotheken gibt, worüber man Ontologien beziehen kann.
Neben Metadaten wie Titel, Lizenz und die neuste Version, werden häufig auch die zugehörigen RDF Daten bereitgestellt.
Die Autoren stellen jedoch fest, dass viele bekannte Ontologie-Bibliotheken nicht mehr existieren bzw. nicht weiterbetrieben werden.
Wir haben die Webseiten der genannten Systeme selbst geprüft und konnten die folgenden weiterhin abrufen:

\begin{enumerate}
    \item BioPortal (\url{https://bioportal.bioontology.org/}, Themengebiete: Biomedizin)
    \item OBO Foundry (\url{https://obofoundry.org/}, Themengebiete: Biologie und Biomedizin)
    \item oeGOV (\url{http://www.oegov.us/}, Themengebiet: e-Government)
    \item Ontology Lookup Service (\url{https://www.ebi.ac.uk/ols4}, Themengebiete: Biomedizin)
    \item Ontology Design Patterns (\url{http://ontologydesignpatterns.org/wiki/Main\_Page}, viele Themengebiete, siehe auch \url{http://ontologydesignpatterns.org/wiki/Community:Domain})
    \item ONKI ontology server (\url{https://onki.fi/en/}, verschiedene Themengebiete)
\end{enumerate}

Nach der Vorstellung vorhandener Ontologie-Bibliohteken stellen verschiedene Charakteristiken dieser Systeme vor.
Daneben gibt es auch Angaben über die konkreten ontologischen Inhalte der Systeme (z.B. Ontology-Mappings, Ontology-Level Relations).
Eine wichtige Erkenntis war, dass die Systeme in der Regel nur einen Ausschnitt an Ontologien anbieten.
Das haben wir selbst bei unserer Recherche festgestellt.
Für die Ermittlung passender Ontologien eines konkreten Fachgebietes ist man gezwungen, mehrere Ontologie-Bibliotheken zu konsultieren, wobei ergänzende Suchen mittels einer Suchmaschine (wie Google) nicht ausbleiben.
Am Ende der Publikation weisen die Autoren noch auf eine Reihe von offenen Problemen solcher Systeme hin.
Als ein Problem wurden die oft mangelhaften Angaben zur Weiterverwendung und Lizensierung von Ontologien genannt.
Ein anderes die oft mangelnde Unterstützung für die Mitarbeit (z.B. Melden von Fehlern, Einreichung von Ergänzungen).
Bei unserer Recherche sind wir auch auf diese Probleme gestoßen.

\subsection{OntoCommons Ontology Catalogue}

https://data.ontocommons.linkeddata.es/index

\subsection{IndustryPortal}

https://industryportal.enit.fr/ontologies

\section{Methodik}

Das IndustryPortal-Projekt und OntoCommons Ontology Catalogue kommen der angestrebten Liste von maschinenlesbaren RDF Ontologien am nächsten.
Jedoch ist der Catalogue nicht ohne zusätzlichen Mehraufwand maschinell zu verarbeiten.
Es handelt es sich um eine Webseite\footnote{Es wird HTML ohne Annotationen der Elemente verwendet, was das Parsing und spätere Auswerten erheblich erschwert.} und bei manchen Ontologien ist kein Link auf RDF-Daten hinterlegt. Abschließend ist noch festzustellen, dass die Liste nur einen kleinen Ausschnitt der verfügbaren Ontologien beinhaltet.

Aus den genannten Gründen wurde eine eigene Literatur- und Internetrecherche durchgeführt und die gefundenen Ontologien in einer CSV-Datei gesammelt.
CSV-Dateien sind einfach zu parsen und werden in jeder gängigen Programmier- und Skriptsprache unterstützt.
Interessenten können die CSV-Datei automatisiert, maschinell auswerten und die referenzierten Ontologien über die angegebenen URLs zu den RDF-Daten direkt konsumieren.
Dies erlaubt durchgängige Maschinenlesbarkeit und einen hohen Grad der Automatisierung.

Die folgenden Forschungsfragen wurden bei der Erstellung der Liste zugrundegelegt:

\begin{enumerate}
    \item Welche Ontologien für das Fachgebiet industrielle Fertigung gibt es?
    \item Welche dieser Ontologien werden aktiv betreut bzw. wann war die letzte dokumentierte Aktivität im Projekt?
\end{enumerate}

Bei der gesamten Recherche waren sowohl deutsch- als auch englischsprachige Inhalte von Interesse.

\subsection{Anlaufstellen für die Suche}

Die Ontologieentwicklung hat starke wissenschaftliche Wurzeln, daher stand die Literaturrecherche am Anfang, um die Publikationen zu finden, welche eine Ontologie und ihre Inhalte näher vorstellen.
Für die Internetsuche wurde \textbf{Google Search} verwendet.
\textbf{Google Scholar}\footnote{\url{https://scholar.google.com/}} kam bei der Recherche nach wissenschaftlichen Publikationen zum Einsatz.
Daneben wurde in weiteren Onlinediensten gesucht, wobei nicht alle unmittelbar etwas mit Ontologien zu tun haben.
Die folgende Liste enthält alle während der Recherche genutzten Onlinedienste:

\begin{enumerate}
    \item \textbf{AURORAL Ontologien} (\url{https://auroral.iot.linkeddata.es/index.html}) - Liste von Ontologien, die für das AURORAL Projekt (\url{https://www.auroral.eu/#/}) entwickelt wurden. Es enthält u.a. Ontologien zum Thema Ladestationen und Energiekonsum.
    \item Projektseite der \textbf{Basic Formal Ontology} (kurz BFO, \url{https://basic-formal-ontology.org/users.html}) - Enthält eine Liste von Ontologie-Projekten, welche die Top-Level Ontologie BFO nutzen.
    \item \textbf{BioPortal} - Ontologie-Portal mit über 1094 Ontologien, vordergründig aus dem biomedizinischen Themenbereich\footnote{Geprüft am 26.03.2024}.
    \item \textbf{Basic Register of Thesauri, Ontologies \& Classifications} (kurz BARTOC, \url{https://bartoc.org}) - Eine Webseite mit angebundener Suchmaschine, die Knowledge Organization Systems, Vokabulare und Ontologien auflistet. Es wurden ausschließlich OWL Ontologien verwendet. Einträge, die nicht in Deutsch oder Englisch vorlagen wurden automatisch übersetzt und geprüft.
    \item \textbf{EU Vocabularies} (\url{https://op.europa.eu/en/web/eu-vocabularies/controlled-vocabularies}) - Eine Webseite u.a. mit Ontologien und Vokabularen, bereitgestellt von der Europäischen Union (bzw. einer ihr unterstellten Institution).
    \item \textbf{Github} (\url{https://www.github.com}) - Ein Onlinedienst zur Softwareentwicklung, der jedoch sehr häufig als Ort für Ontologie-Projekte genutzt wird.
    \item \textbf{IndustryPortal} (\url{https://industryportal.enit.fr/}) - Ontologie-Portal mit über 109 Ontologien\footnote{Geprüft am 26.03.2024.} aus der Industrie und angrenzenden Themenbereichen.
    \item \textbf{Linked Open Vocabularies} (\url{https://lov.linkeddata.es/dataset/lov/}) - Ein kuratierter Katalog von Vokabularen und Ontologien zur Beschreibung von Daten im Internet.
    \item \textbf{OntoCommons Ontology Catalogue} (\url{https://data.ontocommons.linkeddata.es/index}) - Manuell kuratierte Liste mit Ontologien aus den Themenbereichen Industrie, Produktion, Materialwissenschaften, Bauwirtschaften und weiteren.
    \item \textbf{ShowVoc} (\url{https://showvoc.op.europa.eu/}) - Ein Portal mit einer Liste von Ontologien, Vokabularen und anderen. Es wurden nur die Ontologien ausgewertet.
\end{enumerate}

Es wurden die Suchergebnisse für jeden dieser Suchbegriffe näher untersucht, insofern sie von einer Ontologie handelten.

\subsection{Suchparameter}

Eine eindeutige Abgrenzung der industriellen Fertigung zu angrenzenden Themen war mir nicht komplett möglich.
Der Hauptgrund war die Menge an Trends im Rahmen der Industrie 4.0\footnote{\url{https://www.plattform-i40.de/IP/Navigation/EN/Industrie40/WhatIsIndustrie40/what-is-industrie40.html}}.
So spielt die Digitalisierung und starke Vernetzung bei der Produktion eine wichtige Rolle, wodurch die Bedeutung von z.B. Sensoren zunimmt.
Die Bedeutung von Sensoren in der Produktion vor der Industrie 4.0 war im Verhältnis eher untergeordnet.
Im Rahmen der Industrie 4.0 werden Sensoren jedoch häufig von intelligenten Systemen ausgewertet, die aufgrund der Daten Rückschlüsse auf z.B. eine Maschine und deren Einsatz ziehen.
Dafür werden Fakten über die Umwelt verarbeitet, z.B. zu Luftdruck, Gewicht, Materialeigenschaften etc.
Das spielt wiederum bei Digitalen Zwillingen eine wichtige Rolle, welche einen digitalen Repräsentanten einer physischen Entität darstellen.
Sie wiederum erlauben statistische Auswertungen und Simulationen.
Die Liste liese sich weiter fortsetzen.
Anhand dieses Beispiel sollte die Verkettung der unterschiedlichen Themengebiet illustriert werden.

Ontologien befinden sich häufig in Onlinediesten, die eine Schlüsselwort-basierte Suche besitzen.
Daneben gelangt man auch mithilfe von Schlüsselwörtern über Google Search an die Ontologien.
Aus diesem Grund wurde die folgende Liste an Schlüsselwörtern erarbeitet, welche den Themenbereich industrielle Produktion (in einem weiteren Sinne) repräsentieren:

\begin{enumerate}
    \item defect
    \item digital twin
    \item industry
    \item machine
    \item manufacturing
    \item product
    \item supply chain
\end{enumerate}

\textbf{TODO:} Beschreibe die zugehörigen Themenbereiche nochmal... mit Grafik?

Die meisten Onlinedienste besaßen an sich schon wenige Ontologien, weshalb die Suche mit weiteren Schlüsselwörtern oft nicht nötig war, um den gesamten Bestand zu sichten.
Für den Fall, dass es neben Ontologien und Vokabularen noch andere Typen gab, so wurden diese anderen ignoriert.

\subsection{Internetrecherche}

Ergänzt wurde die Literaturrecherche durch eine Internetrecherche, weil nicht alle Ontologien durch wissenschaftliche Literatur dokumentiert sind.
Zur Internetrecherche wurden die folgenden Webseiten genutzt:


TODO: Hin- und Her-Bouncing erklären? Ausgangspunkt > schauen was es an Ontologien gibt > Kriterien verfeinern > ...

\textbf{TODO:}Auf Hilfsskripte eingehen: cleaning CSV file; finde referenzierte Namespaces/Ontologien; liste noch nicht erfasste, aber referenzierte Ontologien aus ontologies.csv auf

\textbf{TODO} Manche Einträge in ontologies.csv wurden inhaltlich angepasst, weil sonst missverständlich, z.B. http://emmo.info/emmo/top/annotations\# hat als Titel "Annotations", bezieht sich aber auf EMMO Ontologie Annotations
- Weiterhin wurde versucht, so wenig wie möglich Abkürzungen zu verwenden

Weiterhin wurden die Import-Statements in den RDF-Daten der Ontologien manuell geprüft und etwaige Verlinkungen in die Recherche aufgenommen.
- Ontologien wie EMMO (http://emmo.info/) haben eine hohe Modularisierung, daher müssen die owl:import Statements genutzt werden, um an alle referenzierten OWL-Daten zu kommen. Die Top- und Mid-Level-Ontologien sind nur "leere" Hüllen.

- Autorenschaft wurde teilweise anhand der Git-Commits ermitteln (z.B. VDI3682), weil direkt in Metadaten und Projektbeschreibung nicht ersichtlich

- die Ontologie RealEstateCore gab ihre OWL-Ontologie in Version 4 auf und verlagerte sich dann auf SHACL und Digital Twin Definition Language Serialisierungen
  - https://dev.realestatecore.io/docs/DTDL-or-SHACL
  - Doc der neusten V3: https://doc.realestatecore.io/3.3/full.html

- BACnet ontology hat 326 prefixes

- RDF Datei (https://doc.realestatecore.io/3.3/full.unioned.rdf) für RealEstateCore war ab 14.03. nicht mehr erreichbar (Page not found error),
  durch kurzes probieren auf full.rdf gestoßen die ging.

- kurios: Product Types Ontology stellt im Dump nur die 1000 am häufigsten abgefragten Klassen als Ontologie bereit (http://www.productontology.org/)

\subsection{Selektion einer Ontologie}

Im Rahmen der Literatur- und Internetrecherche wurde eine Ontologie selektiert, wenn sie die folgenden Kriterien erfüllte:

\begin{enumerate}
    \item Die Ontologie besitzt eine URL zu einer Textdatei, welche in einer formalen Sprache (RDF/OWL) vorliegt.
    \item Die Ontologie stellt ein Vokabular zur Beschreibung des Fachgebietes bereit.
    \item Der Ontologie liegt eine logische Theorie (z.B. in Form von Axiomen, Regeln, Hierarchien) über das Fachgebiet zugrunde, die auf das Vokabular zurückgreift.
    \item Die Inhalte der Ontologie haben einen direkten thematischen Bezug zu den vorher genannten Schlüsselwörtern bzw. den damit zusammenhängenden Themen.
    \item Die Ontologie wurde explizit in einer kuratierten Ontologie-Bibliothek aufgelistet (z.B. OntoCommons ontology catalogue). Es wird in diesem Fall angenommen, dass jemand mit fachlichem Hintergrund die Ontologie bewusst thematisch eingeordnet hat. (TODO)
    \item Enthält eine Instanz von owl:Ontology und x Klassen/Properties definiert
\end{enumerate}

\textbf{TODO: Schlüsselwort => Thema hier angemessen oder sollte das extra behandelt werden?}

Bei jeder selektierten Ontologie wurde anschließend eine Reihe von objektiven Angaben ergänzt.
Eine Bewertung der Inhalte (wie RDF-Daten, genutzte Lizenzen o.ä.) fand nicht statt.
Sie helfen dabei, einen neutralen Blick auf die vorliegenden Industrie-Ontologien zu erhalten.
Im Folgenden die Angaben mit einer kurzen Erklärung:

\begin{itemize}
    \item \textbf{Name der Ontologie} - Der Name der Ontologie, der in den RDF-Daten bzw. der zugehörigen Dokumentation zu finden ist.
    \item \textbf{Ist Industry related?} - TODO
    \item \textbf{Abkürzung} - Falls vorhanden, die Abkürzung des Namens der Ontologie. Wird häufig als Prefix in den RDF-Daten verwendet.
    \item \textbf{Kurzbeschreibung} - Eine kurze, prägnante Beschreibung über den Inhalt der Ontologie.
    \item \textbf{Projektseite oder Publikation} - Eine URL auf die Projektseite, falls vorhanden. Alternativ eine URL zur zugehörigen Publikation.
    \item \textbf{Ontologie IRI} - TODO
    \item \textbf{Neuste Version} - Falls vorhanden, eine Angabe zur neusten Version der Rohdaten. Die neuste Version und das Datum der neusten, dokumentierten Änderung hängen nicht zusammen und können vollkommen verschiedene zeitliche Bezüge haben.
    \item \textbf{Datum der neusten, dokumentierten Änderung} - Eine Datumsangabe der neusten und dokumentierten Änderung an den Rohdaten. Die Änderung kann, muss aber nicht zeitlich mit der neusten Version zusammenfallen.
    \item \textbf{Link zu Rohdaten (RDF/XML oder Turtle-Datei)}: Eine URL zu einer oder mehreren Dateien mit den Rohdaten im RDF-Format.
    \item \textbf{Downloadbereich}: Manche RDF-Dateien besitzen keinen statischen Download-Link oder sind Teil eines ZIP Archivs. Diese findet man dann hier.
    \item \textbf{Autor(en)} - Eine Liste von Namen der Autoren oder beteiligten Gruppen/Unternehmen.
    \item \textbf{Lizenz} - Falls vorhanden, eine Angabe zu der/den verwendeten Lizenz(en).
\end{itemize}

\textbf{TODO:} sofern vorhanden wurde die RDF-Datei im Turtle Format heruntergeladen (leichte Lesbarkeit), ansonsten RDF/XML oder was noch verfügbar war.

\textbf{TODO einarbeiten:} es wurde ein Hilfsskript gebaut, um referenzierte Ontologien/Namespaces zu finden. Dabei sah man, dass auch veraltete Namespaces verwendet werden.
- car options ontology nutzt goodrelations (http://purl.org/goodrelations/v1\#) - Legacy Liste mit diesen Ontologien im PHP-Code (TODO: Liste nennen)

Für nicht RDF-Namespaces wie XML wurde versucht anhand des Kontextes die passende Namespace IRI zu ermitteln

Es gab Ontologien (Information Artifact Ontology aus "Open Biological and Biomedical Ontology Foundry"\footnote{\url{https://obofoundry.org/}}) welche Namespaces auf eine Art verwendeten, deren Zurückverfolgbarkeit unmöglich machte.
So wurde obo:http://purl.obolibrary.org/obo/ verwendet, in den Rohdaten (https://raw.githubusercontent.com/information-artifact-ontology/IAO/master/iao.owl) wurden dann jedoch Einträge wie obo:RO\_0001900 verwendet, welche auf eine Subontology https://obofoundry.org/ontology/ro.html verweisen.
Das ist an sich kein fehlerhaftes Verhalten und mag aus Datenmodellierungssicht Sinn ergeben.
Jedoch erschwert das die Nachverfolgbarkeit und Einordnung der referenzierten Ontologien, wenn statt der eigentlich Ontologie IRI eine IRI auf die Hauptontologie (?) verwendet wird.

- \textbf{HINWEIS für CRAWLING}: Timeout lag bei 5 Sek.
- \textbf{HINWEIS für CRAWLING}: Nur erste 300 Triple der Ontologie eingelesen bzgl. owl:import
- \textbf{HINWEIS für TOOLING}: LibreOffice, VSCode, Docker, PHP (welches Tool für was?)

\subsection{Ignorierte Einträge}

Im Rahmen der Recherche wurden Publikationen zu Ontologien gefunden, bei denen die öffentlich zugängliche Daten nicht ermittelt werden konnten.
Diese Arbeiten wurden aus den folgenden Gründen ignoriert:
\begin{enumerate}
    \item Ohne die zugehörigen Rohdaten in formaler Sprache kann nicht sichergestellt werden, dass eine Ontologie vollständig vorliegt. Davon abgeleitete Arbeiten könnten zu unerwarteten Fehlern und Widersprüchen führen. Weiterhin würde man erzwingen, dass das beschriebenen Wissen zuerst manuell in eine entsprechend Form gebracht werden muss, um eine Weiterverarbeitung zu ermöglichen.
    \item Eine automatisierte Nutzung des modellierten Wissens wird als nicht gegeben angesehen, wenn das Wissen selbst nicht in einer maschinenlesbaren, formalen Sprache vorliegt. Publikationen zu Ontologien liegen häufig als PDF-Datei vor, wobei das darin beschriebene Faktenwissen (inkl. Formeln in Description Logic und Prädikatenlogik 1. Stufe) nicht entsprechend annotiert ist. Für einen Parser sieht es daher genauso aus, wie der restliche Inhalt.
\end{enumerate}


\section{Rechercheergebnisse}

[...]

\textbf{TODO: binde finale Version der Grafik ontology-map.png hier ein}

\textbf{TODO: Zusammenfassung der Rechercherergebnisse}

- Publikationen, welche jedoch auf nicht mehr erreichbare Projektseiten/RDF-Daten verlinken (link rot)

- Es gab Referenzen auf Ontologien, die nicht mehr (in der genutzten Form) existieren (z.B. GoodRelations => gehört nun zu schema.org)

- Manche Ontologien besitzen nur ungenügend Dokumentation zur Herkunft und Autorenschaft: daher Schwierigkeiten herauszufinden, welchen Daten man vertrauen kann

- Manche Ontologien nutzen nicht auflösbare Präfixe, z.B. "Semantically Integrated Planning Model (SIMPM)" nutzt "metadata" > http://data.bioontology.org/metadata/ (gesamte Liste nennen)

- Manche Ontologien haben keine Ontology IRI (z.B. IMAMO)

- Manche Ontologie-Dateien besaßen SSL-Zertifikatsprobleme (alle über industryportal)

- Manche Ontologien besitzen eine IRI, die zu einem 404 Fehler führt, z.B. AMontology > https://github.com/iassouroko/AMontology/blob/master/AMOntology-vFinal1.owl nutzt http://www.semanticweb.org/ina3/ontologies/2016/2/untitled-ontology-122\#
  - oder http://emmo.info/emmo/middle/isq (hier gemeldet: https://github.com/emmo-repo/EMMO/issues/285)

- Nutzung von Namespaces die nicht verwendet wurden und auch auf keine Ontologien zeigten, z.B. Vehicle Sales Ontology zeigt auf http://dbpedia.org/resource/, ohne jedoch nie zu verwenden

- Manchmal wurde die https statt http Version einer ontology IRI verwendet (z.B. Spin Spif => OWL, RDFS)

- Vehicles Sales Ontology nutzt Namespaces wie http://en.wikipedia.org/wiki/, die jedoch weder verwendet werden noch auf gültige Ontologie verweisen

- Sollte es keine RDF-Datei vom Projekt geben, werden auch Mirrors wie DBpedia Archivo genutzt

- Es gab Vorkommen in Portalen, bei denen die Metadaten nicht den Angaben in den RDF-Daten entsprachen, z.B. https://industryportal.enit.fr/ontologies/I40KG
  - Titel im Portal: Industry 4.0 Knowledge Graph
  - Titel in RDF: Industrial IoT Architecture Ontology (DC terms: title property)

Bei der Auswertung der Metadaten wurde die folgende Reihenfolge eingehalten:

\begin{enumerate}
    \item RDF-Daten
    \item In den RDF-Daten verlinkte Projektseite
    \item Angaben auf Onlinedienst, der die Ontologie bereitstellt
\end{enumerate}

Wenn eine Anlaufstelle keine Metadaten zu Tage fördert wurde die nächste probiert.
Und lagen am Ende zu einem Metadatum immer noch keine Angaben vor, so wurde an der Stelle "Information not available" in der CSV-Datei hinterlegt.

%----------------------------------------
% eingeschlossenen Textblock kürzen
%

  - Beispiel Volkswagen Vehicles Ontology: Aufgrund fehlender Originalquelle (z.B. wiss. Publikation oder Projektseite) und der unterschiedlichen Dokumentation zur Volkswagen Vehicles Ontology, werden die aktuellen Funde im Folgenden einzeln dargelegt.

Die Ontologie wird in einer Untersuchung namens "Case Study: Contextual Search for Volkswagen and the Automotive Industry" des W3C aus dem Jahre 2011 erwähnt (\cite{greenly2011case}, S. 2).
In der Studie ging es u.a. um die Verbesserung der Suche für Nutzer in digitalen Quellen.
In diesem Zuge wurde u.a. die Volkswagen Vehicles Ontology entwickelt, um Volkswagen-spezifische Konzepte bei Fahrzeugen beschreiben zu können.
Der Autor der Ontologie war Martin Hepp.
Man nutzte diese und andere Vokabulare/Ontologien, um die eigenen Daten anzureichern, um kontextbasierte Suchanfragen zu ermöglichen.
Die URL http://purl.org/vvo/ns wird als Quelle für die Ontologie angegeben, jedoch gelangt man darüber am Ende nur zu https://www.volkswagen.co.uk/en.html, welches keine weiteren Informationen bereitstellt.
Weitere Details zur Ontologie selbst wurden nicht genannt.

Der Link https://industryportal.enit.fr/ontologies/VVO führt zu einem Projekteintrag auf IndustryPortal.
Als Autor wird Abdelouadoud Rasmi angegeben, der jedoch keine ersichtliche Verbindung zu Volkswagen hat.
In der zugehörigen OWL-Datei\footnote{https://data.industryportal.enit.fr/ontologies/VVO/submissions/1/downloadapikey=019adb70-1d64-41b7-8f6e-8f7e5eb54942} findet man jedoch im Abschnitt Autor auch Martin Hepp und bei Beitragende (contributors) u.a. den Erstautor der oben genannten W3C Studie William Greenly.
Aufgrund dieses und anderer Gemeinsamkeiten, z.B. gleicher Name (Volkswagen Vehicles Ontology) und XML-Namensraum (http://purl.org/vvo/ns\#), wird angenommen, dass es sich um die gleiche Ontologie handelt. Die Ontologie auf IndustryPortal hat als Veröffentlichungsdatum den 2. November 2022, wobei als Veröffentlichungsdatum bei der Version der 12. Oktober 2010 genannt wird.
Es liegen keine Informationen vor, dass zwischen diesen Daten weitere Veränderungen an der Ontologie vorgenommen wurden.
Ich gehe daher davon aus, dass es sich hier lediglich um einen erneuten Upload der Version von 2010 handelt.

%----------------------------------------



Weiterführende Metadaten zu den Ontologien finden Sie in unseren Rohdaten, welche in Form einer CSV-Datei öffentlich zur Verfügung stehen. \textbf{TODO Link zu Github Repository}

Im Anhang ist eine Tabelle aller selektierten Ontologien in Form einer Tabelle zu finden.


\section{Fazit und weitere Arbeiten}

Vorschläge für Verbesserungen?

[...]

\subsection{Erwähnungen}

Ich danke Paul-Robert Kästner für die Sichtung und Zuarbeit der ersten Ontologie-Einträge.

Die Arbeit wurde im Rahmen KI-Werk Projektes durchgeführt, welches vom Bundesministerium für Bildung und Forschung (BMBF) gefördert wurde (https://www.cbasynergy.net/cba/ki-werk.html).



\medskip

\printbibliography

\end{document}

